% results.tex
\section{Results}
\vspace{-0.33cm}


\subsection{Representative Shapes and Key Results}
We focus on Llama-style decode GEMMs with small batch ($M \le 8$) and large hidden
dimensions ($K,N \in \{4096,14336\}$), plus a narrow KV projection ($N=1024$).
Wall-clock timings are from our benchmark harness (profilers replay kernels and
inflate timings); these shapes dominate autoregressive decode in Llama-style
transformers.

% Keep your existing Table 1 (representative shapes) wherever it currently lives.
% Keep your existing summary_table.tex input wherever it currently lives.
\begin{table*}[t]
  \centering
  \caption{Summary of key decode shapes (M=1).}
  \label{tab:summary}
  \begin{tabular}{lccccl}
    \toprule
    Shape & FP16 (ms) & INT4 (ms) & Speedup & Bound \\
    \midrule
    KV proj & 0.027 & 0.043 & 0.62x & Dequant \\
    Q proj & 0.075 & 0.047 & 1.58x & Mixed \\
    FFN up & 0.239 & 0.067 & 3.58x & Memory \\
    FFN down & 0.258 & 0.152 & 1.69x & Memory \\
    \bottomrule
  \end{tabular}
\end{table*}


\subsection{Where Does INT4 Help?}
We first identify the empirical performance regimes across representative
decode GEMMs. Figure~\ref{fig:speedup-vs-n} reports speedup relative to FP16
as output width $N$ increases at fixed $K=4096$, isolating the effect of layer
dimension on performance.

\begin{figure}[h]
\centering
\begin{subfigure}{\linewidth}
  \includegraphics[width=\linewidth]{speedup_vs_n_k4096.png}
  \vspace{-0.7cm}
  \caption{Speedup vs output width.}
  \label{fig:speedup-vs-n}
\end{subfigure}

\vspace{-0.05cm}

\begin{subfigure}{\linewidth}
  \includegraphics[width=\linewidth]{family_latency_m1.png}
  \vspace{-0.7cm}
  \caption{Latency by layer family.}
  \label{fig:family-latency-m1}
\end{subfigure}
\vspace{-0.7cm}

\caption{Decode performance across representative shapes.}
%\label{fig:decode-performance}
\end{figure}

INT4 underperforms for narrow projections (e.g., KV projections with small $N$),
but delivers substantial gains as matrix width increases, reaching up to
$3.7\times$ speedup for wide FFN layers. Figure~\ref{fig:family-latency-m1}
confirms this behavior at the layer-family level: projection layers show limited
or negative benefit, while FFN layers consistently improve. The KV projection
regression persists for both $M=1$ and $M=8$, indicating a geometry-driven
transition rather than a batch-size effect. Across all $M \in \{1,\dots,64\}$,
the transition remains near $N \approx 2$--$4$K for $K=4096$.

\subsection{Prefill vs Decode}
We compare prefill-like ($M \gg 1$) and decode-like ($M \le 8$) regimes using
the same layer-family grouping (Figure~\ref{fig:prefill-decode}). Prefill
increases effective arithmetic intensity by amortizing quantization overhead
across more work, yielding consistent INT4 speedups across shapes.

Decode operates in a latency-dominated regime where small batch sizes expose
fixed overheads such as dequantization and launch cost, making INT4 benefits
shape-dependent. For example, FFN-up improves by roughly 3.5$\times$ in decode
while KV projection remains below 1$\times$; in prefill, FFN-up is about 3.0$\times$
and KV improves to roughly 2.3$\times$.

\vspace{0.25cm}
\begin{figure}[h]
  \centering
  \includegraphics[width=\linewidth]{prefill_vs_decode_speedup.png}
  \vspace{-0.6em}
  \caption{Prefill vs decode speedup by family.}
  \label{fig:prefill-decode}
  \vspace{-0.8em}
\end{figure}

\subsection{Why Does INT4 Help (or Hurt)?}
We attribute the observed performance regimes using hardware counters and a
targeted dequantization microbenchmark. Together, these measurements explain
why INT4 provides large gains for some decode shapes while regressing on others.

Figure~\ref{fig:util-compare} compares hardware utilization between FP16 and
INT4 kernels. On the A10G, FP16 decode GEMMs operate close to the bandwidth
limit: memory throughput reaches approximately 75--77\% of peak DRAM bandwidth
while achieved compute throughput remains relatively low, indicating a strongly
memory-bound regime where performance is limited by weight movement rather than
arithmetic capability. INT4 weight packing reduces memory traffic and shifts the execution regime.

\begin{figure}[h]
  \centering
  \includegraphics[width=\linewidth]{ncu_utilization_comparison.png}
  \vspace{-0.6em}
  \caption{Hardware utilization comparison (FP16 vs INT4).}
  \label{fig:util-compare}
\end{figure}

However, reduced memory traffic alone does not guarantee speedup. The
dequantization microbenchmark in Figure~\ref{fig:dequant-breakdown}
isolates the cost of unpacking INT4 weights during execution. For narrow
projection layers (e.g., KV projections), total work per launch is small,
and the fixed dequantization overhead constitutes a significant fraction
of runtime. As a result, INT4 kernels can underperform FP16 despite lower
memory traffic.

In contrast, wide FFN projections provide sufficient arithmetic work to
amortize dequantization cost, allowing INT4 to achieve substantial latency
improvements. Overall, INT4 improves decode performance when reduced memory
movement outweighs added dequantization overhead.

\begin{figure}[h]
  \centering
  \includegraphics[width=\linewidth]{dequant_breakdown.png}
  \vspace{-0.6em}
  \caption{Dequantization breakdown for representative decode shapes.}
  \label{fig:dequant-breakdown}
  \vspace{-0.8em}
\end{figure}

\subsection{A Regime Model for INT4 Decode GEMMs}
We summarize the observed performance behavior using a simple latency
decomposition:
\begin{equation}
T_{\text{total}} = T_{\text{launch}} + T_{\text{mem}}(W) +
T_{\text{dequant}} + T_{\text{compute}}.
\end{equation}

This formulation captures the competing costs exposed in decode workloads.
Narrow projections perform limited arithmetic per launch and are therefore
dominated by fixed overheads, particularly $T_{\text{dequant}}$ and launch
latency. In contrast, wide FFN layers perform substantially more work per
memory access, making execution primarily bandwidth-limited through
$T_{\text{mem}}$.

INT4 weight packing reduces memory traffic, decreasing $T_{\text{mem}}$ and
shifting execution toward a mixed compute--memory regime. Compute utilization
and roofline views support this model. INT4 achieves \~23\% peak SM throughput
vs \~32\% for FP16 (about 28\% lower), reflecting reduced compute pressure under
the INT4 kernel. The roofline view in Figure~\ref{fig:roofline} illustrates this
shift away from bandwidth saturation.

\begin{figure}[h]
  \centering

  \begin{subfigure}[t]{0.4\linewidth}
    \centering
    \includegraphics[width=\linewidth]{peak_compute_utilization.png}
    \caption{Peak compute utilization (SM throughput).}
    \label{fig:peak-compute}
  \end{subfigure}
  \hfill
  \begin{subfigure}[t]{0.58\linewidth}
    \centering
    \includegraphics[width=\linewidth]{roofline_scatter.png}
    \caption{Roofline view (arithmetic intensity vs TFLOPs).}
    \label{fig:roofline}
  \end{subfigure}

  \vspace{-0.6em}
  \caption{Utilization regime view. INT4 shifts execution away from bandwidth
  saturation and toward a mixed compute--memory regime.}
  \label{fig:regime-utilization}
\end{figure}

INT4 weight packing introduces a transition boundary in arithmetic-intensity
space separating overhead-limited and bandwidth-limited regimes. In our sweep,
the boundary occurs near $\alpha \approx 8$ FLOPs/byte.

\begin{figure}[h]
  \centering
  \includegraphics[width=\linewidth]{regime_boundary.png}
  \vspace{-0.8cm}
  \caption{Regime boundary: speedup vs arithmetic intensity.}
  \label{fig:regime-boundary}
\end{figure}

Figure~\ref{fig:regime-boundary} visualizes this transition. Shapes below the
boundary remain dominated by fixed overheads and therefore see limited benefit
from INT4, while shapes above it achieve consistent speedups as reduced memory
movement outweighs dequantization cost.

This model consolidates the empirical observations from earlier sections:
narrow projections operate below the boundary, while increasing $N$ raises
arithmetic intensity and shifts execution into a bandwidth-dominated regime
where INT4 becomes advantageous.

\begin{center}
\fbox{\parbox{0.97\linewidth}{
\textbf{INT4 Decision Rule.} INT4 improves decode latency when the memory savings
from FP16 $\rightarrow$ INT4 exceed dequantization cost:
$T_{\text{mem}}^{\text{FP16}} - T_{\text{mem}}^{\text{INT4}} > T_{\text{dequant}}$.
}}
\end{center}

In practice, this boundary provides a deployment guideline:
quantization is most effective for wide decode GEMMs and FFN layers,
while narrow projection layers may benefit more from kernel fusion or
launch amortization strategies than from aggressive weight compression.

\subsection{Systems-Level View}
Because autoregressive decoding executes one token at a time, we connect kernel-level behavior to end-to-end decode constraints through batch scaling and throughput analysis. These latency–throughput tradeoffs directly determine user-visible response time in interactive LLM serving.

Figure~\ref{fig:batch-scaling-and-tps} shows that latency grows sublinearly
with batch size while throughput increases substantially, reflecting improved
hardware utilization as launch overhead is amortized. 

\begin{figure}[h]
  \centering

  \begin{subfigure}{\linewidth}
    \centering
    \includegraphics[width=\linewidth]{batch_scaling_k4096_n4096.png}
    \vspace{-0.6cm}
    \caption{Latency vs batch ($K=N=4096$).}
    \label{fig:batch-scaling}
  \end{subfigure}

  %\vspace{-0.6em}

  \begin{subfigure}{\linewidth}
    \centering
    \includegraphics[width=\linewidth]{tokens_per_sec_k4096_n4096.png}
    \vspace{-0.5cm}
    \caption{Throughput (tokens/sec) vs batch ($K=N=4096$).}
    \label{fig:tps}
  \end{subfigure}

  \vspace{-0.5em}
  \caption{Batch scaling tradeoffs for a representative decode shape.}
  \label{fig:batch-scaling-and-tps}
  \vspace{-0.8em}
\end{figure}

  \vspace{0.8em}

For small batches ($M=1$–$4$), decode latency is dominated by fixed kernel overheads, limiting the benefit of additional compute efficiency. As batch size increases, INT4 benefits become more consistent because memory traffic and compute utilization rather than launch overhead determine performance.

These results highlight an important systems tradeoff: optimizations that
improve single-token latency do not always maximize throughput, and deployment
targets (interactive latency vs. serving throughput) determine whether INT4
provides meaningful gains. This behavior mirrors production serving systems, where batching improves throughput but increases token latency, creating a practical optimization tradeoff.


\subsection{Regime Map and Kernel Composition}
We connect kernel-level behavior to system-level impact by examining where
decode runtime is actually spent and how performance regimes distribute
across model shapes.

Figure~\ref{fig:kernel-time} shows that a small set of wide FFN GEMMs dominates
total CUDA execution time during decode. Although transformer layers contain
many projection operations, runtime is heavily concentrated in
high-dimensional FFN computations. Consequently, optimizing these shapes yields
disproportionate end-to-end latency improvements, motivating decode-focused
specialization rather than uniform optimization across all layers.

\begin{figure}[h]
  \centering
  \includegraphics[width=\linewidth]{kernel_time_top_shapes.png}
  \vspace{-0.6cm}
  \caption{Top shapes by CUDA time (Nsight Systems).}
  \label{fig:kernel-time}
\end{figure}

Figure~\ref{fig:heatmap} complements this view by visualizing performance
regimes across decode shapes. INT4 provides consistent speedups only once
arithmetic intensity exceeds a threshold determined primarily by output width
and hidden dimension. Regions with small $N$ remain dominated by
dequantization and launch overhead, explaining regressions observed for
KV-like projections.


\begin{figure}[h]
  \centering
  \includegraphics[width=\linewidth]{int4_speedup_heatmap_m1.png}
  \vspace{-0.8cm}
  \caption{Decode regime map (M=1) across $K$ and $N$.}
  \label{fig:heatmap}
\end{figure}

Together, these results unify the regime analysis with runtime behavior: the
layers that dominate decode execution are those above the regime boundary where
INT4 provides the largest benefit.
